\documentclass[10pt,a4paper]{article}

%% paquetes
\usepackage[utf8]{inputenc}
\usepackage{a4wide,color,verbatim,Sweave,url,xargs,amsmath,booktabs,longtable,eurosym}
\usepackage{tikz}
\usetikzlibrary{automata,positioning}

%% nuevos entornos
\newenvironment{question}{\item}{}
\newenvironment{solution}{\comment}{\endcomment}
\newenvironment{answerlist}{\renewcommand{\labelenumii}{(\alph{enumii})}\begin{enumerate}}{\end{enumerate}}

%% comandos para exams
\newcommand{\exname}[1]{\def\@exname{#1}}
\newcommand{\extype}[1]{\def\@extype{#1}}
\newcommand{\exsolution}[1]{\def\@exsolution{#1}}
\newcommand{\extol}[1]{\def\@extol{#1}}
\newcommand{\exclozetype}[1]{\def\@exclozetype{#1}}
\newcommand{\exshuffle}[1]{\def\@exshuffle{#1}}

%% párrafos
\setlength{\parskip}{0.7ex plus0.1ex minus0.1ex}
\setlength{\parindent}{0em}

\begin{document}
\Sconcordance{concordance:automaton.tex:automaton.Rnw:1 29 1 1 56 5 1 1 6 3 1 1 2 5 0 %
1 2 9 1 1 2 5 0 1 2 9 1}


\begin{enumerate}



\begin{question}

Considere el siguiente autómata con estado inicial A y estado de aceptación A:

\includegraphics[width=5cm]{automaton.png}

¿Cuáles de las siguientes secuencias son aceptadas?

\begin{answerlist}
  \item 1011
  \item 11110
  \item 10110
  \item 010110
  \item 0101
\end{answerlist}
\end{question}


\begin{solution}

El autómata dado acepta cadenas de entrada que consisten en un número
par de unos y un número
par de ceros.

\begin{answerlist}
  \item No aceptada
  \item No aceptada
  \item No aceptada
  \item No aceptada
  \item Aceptada
\end{answerlist}
\end{solution}


\exname{Autómata}
\extype{mchoice}
\exsolution{00001}

\end{enumerate}
\end{document}
